\documentclass[10pt,a4paper]{amsart}
\usepackage[margin=19mm]{geometry}
\usepackage{amsmath,amssymb,mathtools}
\usepackage[scr=rsfs,cal=euler]{mathalfa}
\usepackage{xfrac}
\usepackage[unicode,hidelinks,bookmarks=false]{hyperref}
\newcommand{\ie}{i.e.\ }
\newcommand{\cf}{cf.\ }
\newcommand{\resp}{respectively\ }
\newcommand{\Subsection}{\subsection}
\providecommand{\nobreakdash}{\nobreak\hbox{-}}
\setlength{\emergencystretch}{2em}
\multlinegap0pt
\usepackage{amssymb}
\usepackage{mathtools}
\usepackage[shortlabels]{enumitem}
\newtheorem{lemma}{Lemma}
\title{Multicolor Ramsey and list Ramsey numbers for star-like trees}
\author{Qinghong Zhao, Yaping Mao, Xiangqian Zhou}
\date{Source: arXiv:2609.06378v1 (CC BY 4.0)}
\begin{document}
\maketitle

\noindent\emph{Source and licence.} Multicolor Ramsey and list Ramsey numbers for star-like trees. Version arXiv:2609.06378v1 (CC BY 4.0), distributed by arXiv under the Creative Commons Attribution 4.0 International licence, http://creativecommons.org/licenses/by/4.0/, as recorded in the arXiv licence metadata for this exact version and shown on the abstract page https://arxiv.org/abs/2609.06378v1. Full licence notice: \texttt{licenses/arxiv-2609.06378-CC-BY-4.0.txt}.\par\smallskip

\noindent\textit{Editorial scope.} Counted unit: the article's lemma lem:extension (source0.tex lines 627--643). The article's complete original proof of it is delivered with it (source0.tex lines 645--704, one \texttt{proof} environment of 60 lines). No mathematics is written, repaired or re-derived: the statement and the proof are the article's own lines, in the article's own notation, and no retained line is substituted or re-flowed.  No further source-local context is needed: every cross-reference of every retained line resolves inside the two delivered spans.  Nothing was omitted from any retained span, so the package carries no omission record.  No retained line contains a citation, so the wrapper carries no bibliography and no bibliography entry.  No retained line contains a figure environment, an image-inclusion command or drawing code, so no image is delivered and the package declares no asset dependency.  Editorial formatting only, in addition: an AMS \texttt{amsart} preamble that restates the article's own declarations and macros used by the retained lines, namely the environments \texttt{lemma} on separate counters, together with the standard packages those lines need.  The retained environments print in this document as Lemma 1 in order of appearance, whereas the article prints the same items with the numbering of its own full document.  The complete native source of the article as distributed in the arXiv e-print for this exact version is bundled unchanged as \texttt{sources/209530/source0.tex}.
\begin{lemma}\label{lem:extension}
Let \(n,q\) be integers with \(q\geq3\), let
\(\varepsilon,\eta\in\{0,1\}\), and suppose that
\(n\geq5q-7+2\varepsilon\).
Let \(G\) be a red--blue complete graph, and let \(A\) and \(B\) be
disjoint subsets of \(V(G)\) such that
\(|A|=n-q+1\), \(|B|=n-q-\varepsilon\),
and
\(\Delta(G_B[A])\leq q+\varepsilon-1\), and \(\Delta(G_B[B])\leq q-3\).
Let \(c\notin A\cup B\) and suppose that
\(d_R(c,A\cup B)\geq n-\eta\).
Suppose further that there are distinct vertices
\(u,v\notin A\cup B\cup\{c\}\) such that \(cuv\) is a red path. If
\(\eta=1\), suppose in addition that there is a vertex
\(w\notin A\cup B\cup\{c,u,v\}\) such that \(cw\) is red. Then
\(G\) contains a red copy of \(S_n^q\) centered at \(c\).
\end{lemma}

\begin{proof}
For \(X\in\{A,B\}\), let \(N_X:=N_{G_R}(c)\cap X\) and
\(M_X:=X\setminus N_X\), and define \(r_X:=|N_X|\),
\(t_X:=|M_X|\), \(D_A:=q+\varepsilon-1\), and \(D_B:=q-3\).
If \(\mu_X\) denotes the maximum size of a red matching between
\(N_X\) and \(M_X\), then
\(\mu_X\geq\min\{r_X,(t_X-D_X)_+\}\),
where \((z)_+:=\max\{z,0\}\). Indeed, if a maximum such matching
does not saturate \(N_X\), an uncovered vertex of \(N_X\) has at
least \((t_X-D_X)_+\) red neighbors in \(M_X\), all of which are
covered by the matching.

Let \(d:=r_A+r_B\), \(a:=\min\{q-1,d-n+1+\eta\}\), and
\(b:=q-1-a\).
Since \(d\geq n-\eta\), we have \(a\geq1\) and \(b\leq q-2\).
Moreover,
\(r_A\geq d-|B|\geq q+\varepsilon-\eta\geq b\) and \(r_B\geq d-|A|\geq q-1-\eta\geq b\).
Assume first that \(b>0\). Then \(a=d-n+1+\eta\), and hence
\(b=n+q-2-\eta-d\).
Therefore,
\begin{align*}
 (t_A-D_A)_++(t_B-D_B)_+
 &\geq t_A+t_B-D_A-D_B\\
 &=2n-4q+5-2\varepsilon-d\\
 &\geq n+q-2-\eta-d=b.
\end{align*}
Let \(p_X:=(t_X-D_X)_+\) for \(X\in\{A,B\}\).
Since \(r_A,r_B\geq b\), the preceding matching bound and
\(p_A+p_B\geq b\) imply
\(\mu_A+\mu_B\geq \min\{b,p_A\}+\min\{b,p_B\} \geq \min\{b,p_A+p_B\}=b\).
For each \(X\in\{A,B\}\), choose a red matching \(F_X\) of size
\(\mu_X\) between \(N_X\) and \(M_X\). Since \(A\cap B=\varnothing\),
\(F_A\cup F_B\) is a matching of size at least \(b\); let \(F\) be
any \(b\)-edge submatching of \(F_A\cup F_B\). If \(b=0\), let
\(F=\varnothing\).

Delete from \(N_A\cup N_B\) the \(b\) endpoints of \(F\) lying in
that set, and denote the remaining set by \(N'\). We claim that
\(G_R[N']\) contains a matching of size \(a\). Otherwise, deleting
the endpoints of a maximum matching leaves a red-independent set
whose intersections with \(A\) and \(B\) are blue cliques of orders
at most \(D_A+1\) and \(D_B+1\), respectively. Therefore,
\(d-b-2(a-1)\leq D_A+D_B+2\).
Since \(a+b=q-1\), it follows that
\(d-a\leq D_A+D_B+q-1=3q+\varepsilon-5\).
On the other hand,
\(d-a =\max\{d-q+1,n-1-\eta\} \geq n-1-\eta >3q+\varepsilon-5\),
where the strict inequality follows from
\(n\geq5q-7+2\varepsilon\), \(q\geq3\), and \(\eta\leq1\).
This is a contradiction.

The matching \(F\) and a red matching of size \(a\) in \(N'\)
form \(a+b=q-1\) red paths of length two beginning at \(c\) that are
pairwise vertex-disjoint outside \(c\). They use \(b+2a\) red neighbors of \(c\) in
\(A\cup B\), leaving at least
\(d-b-2a=d-q+1-a\geq n-q-\eta\) such neighbors unused. These paths,
the path \(cuv\), the edge
\(cw\) when \(\eta=1\), and \(n-q-\eta\) unused red neighbors of
\(c\) form a red copy of \(S_n^q\) centered at \(c\).
\end{proof}

\end{document}
